\chapter{2011年北京卷（文）}

\section{单选题}

\begin{problem}
已知全集 \(U = \R\)，集合 \(P = \{x \mid x^2 \leqslant 1\}\)，那么 \(\complement_U P =\)（\quad）

\choices
    {\((-\infty, -1]\)}
    {\([1, +\infty)\)}
    {\((-1, 1)\)}
    {\((-\infty, -1) \cup (1, +\infty)\)}
\end{problem}

\begin{answer}
D
\end{answer}

\begin{solution}
由 \(x^{2}\leqslant 1\)，得 \(-1\leqslant x\leqslant 1\)，所以
\[
P=[-1,1]
\]
全集为 \(\R\)，因此
\[
\complement_U P=(-\infty,-1)\cup(1,+\infty)
\]
故选 D.
\end{solution}

\begin{problem}
复数 \(\frac{\i - 2}{1 + 2\i} =\)（\quad）

\choices
    {\(\i\)}
    {\(-\i\)}
    {\(-\frac{4}{5} - \frac{3}{5}\i\)}
    {\(-\frac{4}{5} + \frac{3}{5}\i\)}
\end{problem}

\begin{answer}
A
\end{answer}

\begin{solution}
分子分母同乘分母的共轭复数 \(1-2\i\)，得
\[
\frac{\i-2}{1+2\i}
=\frac{(\i-2)(1-2\i)}{(1+2\i)(1-2\i)}
=\frac{5\i}{5}
=\i
\]
故选 A.
\end{solution}

\begin{problem}
如果 \(\log_{\frac{1}{2}} x < \log_{\frac{1}{2}} y < 0\)，那么（\quad）

\choices
    {\(y < x < 1\)}
    {\(x < y < 1\)}
    {\(1 < x < y\)}
    {\(1 < y < x\)}
\end{problem}

\begin{answer}
D
\end{answer}

\begin{solution}
因为底数 \(\frac{1}{2}\in(0,1)\)，函数
\(\log_{\frac{1}{2}}x\) 在 \((0,+\infty)\) 上单调递减.
由
\[
\log_{\frac{1}{2}}x<\log_{\frac{1}{2}}y
\]
可得 \(x>y\). 又由
\[
\log_{\frac{1}{2}}y<0=\log_{\frac{1}{2}}1
\]
可得 \(y>1\). 因而 \(1<y<x\)，故选 D.
\end{solution}

\begin{problem}
若 \(p\) 是真命题，\(q\) 是假命题，则（\quad）

\choices
    {\(p \wedge q\) 是真命题}
    {\(p \vee q\) 是假命题}
    {\(\neg p\) 是真命题}
    {\(\neg q\) 是真命题}
\end{problem}

\begin{answer}
D
\end{answer}

\begin{solution}
\(p\) 为真命题、\(q\) 为假命题，所以 \(p\wedge q\) 为假命题，
\(p\vee q\) 为真命题，\(\neg p\) 为假命题，\(\neg q\) 为真命题.
因此只有选项 D 正确.
\end{solution}

\begin{problem}
某四棱锥的三视图如图所示，该四棱锥的表面积是（\quad）

\texfigure[max width=0.58\linewidth]{img_repaint/2011/beijing_liberal/q5_fig1.tex}

\choices
    {\(32\)}
    {\(16 + 16\sqrt{2}\)}
    {\(48\)}
    {\(16 + 32\sqrt{2}\)}
\end{problem}

\begin{answer}
B
\end{answer}

\begin{solution}
由三视图可知，该四棱锥的底面为边长 \(4\) 的正方形，顶点在底面
中心的正上方，棱锥的高为 \(2\). 因而底面面积为 \(16\).
从底面中心到任一边的距离为 \(2\)，所以每个侧面的高为
\[
\sqrt{2^{2}+2^{2}}=2\sqrt{2}
\]
每个侧面是底边为 \(4\) 的三角形，面积为
\[
\frac{1}{2}\times4\times2\sqrt{2}=4\sqrt{2}
\]
四个侧面面积之和为 \(16\sqrt{2}\)，故四棱锥的表面积为
\[
16+16\sqrt{2}
\]
故选 B.
\end{solution}

\begin{problem}
执行如图所示的程序框图，若输入 \(A\) 的值为 2，则输出的 \(P\) 的值为（\quad）

\texfigure[max width=0.38\linewidth]{img_repaint/2011/beijing_liberal/q6_fig1.tex}

\choices
    {\(2\)}
    {\(3\)}
    {\(4\)}
    {\(5\)}
\end{problem}

\begin{answer}
C
\end{answer}

\begin{solution}
程序开始时 \(P=1\)，\(S=1\). 每次循环先令 \(P\) 增加 \(1\)，
再令 \(S\) 加上新的 \(P\) 的倒数，循环条件为 \(S\leqslant A=2\).
第一次循环后
\(P=2\)，\(S=1+\frac{1}{2}=\frac{3}{2}\).
第二次循环后
\(P=3\)，\(S=\frac{3}{2}+\frac{1}{3}=\frac{11}{6}\).
第三次循环后
\(P=4\)，\(S=\frac{11}{6}+\frac{1}{4}=\frac{25}{12}>2\).
此时停止循环并输出 \(P=4\)，故选 C.
\end{solution}

\begin{problem}
某车间分批生产某种产品，每批的生产准备费用为 \(800\text{ 元}\). 若每批生产 \(x\) 件，则平均仓储时间为 \(\frac{x}{8}\text{ 天}\)，且每件产品每天的仓储费用为 \(1\text{ 元}\). 为使平均到每件产品的生产准备费用与仓储费用之和最小，每批应生产产品（\quad）

\choices
    {\(60\text{ 件}\)}
    {\(80\text{ 件}\)}
    {\(100\text{ 件}\)}
    {\(120\text{ 件}\)}
\end{problem}

\begin{answer}
B
\end{answer}

\begin{solution}
每件产品平均分摊的生产准备费用为 \(\frac{800}{x}\) 元，
平均仓储费用为 \(\frac{x}{8}\) 元，所以总费用为
\(g(x)=\frac{800}{x}+\frac{x}{8}\)，\(x>0\).
由基本不等式，
\[
g(x)\geqslant 2\sqrt{\frac{800}{x}\cdot\frac{x}{8}}
=20
\]
等号成立当且仅当
\[
\frac{800}{x}=\frac{x}{8}
\]
即 \(x^{2}=6400\). 由于 \(x>0\)，得 \(x=80\).
故选 B.
\end{solution}

\begin{problem}
已知点 \(A(0,2)\)，\(B(2,0)\). 若点 \(C\) 在函数 \(y = x^2\) 的图象上，则使得 \(\triangle ABC\) 的面积为 2 的点 \(C\) 的个数为（\quad）

\choices
    {\(4\)}
    {\(3\)}
    {\(2\)}
    {\(1\)}
\end{problem}

\begin{answer}
A
\end{answer}

\begin{solution}
设 \(C=(u,u^{2})\). 直线 \(AB\) 的方程为
\[
x+y-2=0
\]
且 \(AB=2\sqrt{2}\). 因此
\[
[\triangle ABC]
=\frac{1}{2}\cdot2\sqrt{2}\cdot
\frac{|u+u^{2}-2|}{\sqrt{2}}
=|u^{2}+u-2|
\]
令该面积等于 \(2\)，得到
\[
|u^{2}+u-2|=2
\]
于是
\[
u^{2}+u-2=2
\quad\text{或}\quad
u^{2}+u-2=-2
\]
前一个方程为 \(u^{2}+u-4=0\)，判别式为 \(17>0\)，有两个实根；
后一个方程为 \(u^{2}+u=0\)，有两个实根 \(u=0\)，\(-1\).
四个根彼此不同，故共有 \(4\) 个点，选 A.
\end{solution}

\section{填空题}

\begin{problem}
在 \(\triangle ABC\) 中，若 \(b = 5\)，\(\angle B = \frac{\pi}{4}\)，\(\sin A = \frac{1}{3}\)，则 \(a =\)\fillinblank{}.
\end{problem}

\begin{answer}
\(\frac{5\sqrt{2}}{3}\)
\end{answer}

\begin{solution}
由正弦定理，
\[
\frac{a}{\sin A}=\frac{b}{\sin B}
\]
代入 \(b=5\)、\(\sin A=\frac{1}{3}\)、\(\sin B=\frac{\sqrt{2}}{2}\)，得
\[
a=\frac{5\sin A}{\sin B}
=\frac{5\cdot\frac{1}{3}}{\frac{\sqrt{2}}{2}}
=\frac{5\sqrt{2}}{3}
\]
\end{solution}

\begin{problem}
已知双曲线 \(x^{2}-\frac{y^{2}}{b^{2}}=1\)（\(b>0\)）的一条渐近线的方程为 \(y = 2x\)，则 \(b=\)\fillinblank{}.
\end{problem}

\begin{answer}
\(2\)
\end{answer}

\begin{solution}
将双曲线方程中的常数项 \(1\) 换为 \(0\)，得到渐近线方程
\[
x^{2}-\frac{y^{2}}{b^{2}}=0
\]
即
\[
y=\pm bx
\]
已知一条渐近线为 \(y=2x\)，且 \(b>0\)，所以 \(b=2\).
\end{solution}

\begin{problem}
已知向量 \(\bs{a} = (\sqrt{3}, 1)\)，\(\bs{b} = (0, -1)\)，\(\bs{c} = (k, \sqrt{3})\). 若 \(\bs{a} - 2\bs{b}\) 与 \(\bs{c}\) 共线，则 \(k =\)\fillinblank{}.
\end{problem}

\begin{answer}
\(1\)
\end{answer}

\begin{solution}
\[
\bs{a}-2\bs{b}=(\sqrt{3},1)-2(0,-1)=(\sqrt{3},3)
\]
两向量共线，所以它们的行列式为零：
\[
\begin{vmatrix}
\sqrt{3} & 3\\
k & \sqrt{3}
\end{vmatrix}
=3-3k=0
\]
因此 \(k=1\).
\end{solution}

\begin{problem}
在等比数列 \(\left\{a_{n}\right\}\) 中，\(a_{1}=\frac{1}{2}\)，\(a_{4}=-4\)，则公比 \(q=\)\fillinblank{}；\(a_{1}+a_{2}+\cdots+a_{n}=\)\fillinblank{}.
\end{problem}

\begin{answer}
\(-2\)；\(\frac{1-(-2)^{n}}{6}\)
\end{answer}

\begin{solution}
由等比数列通项公式，
\[
a_{4}=a_{1}q^{3}
\]
所以
\(-4=\frac{1}{2}q^{3}\)，\(q^{3}=-8\)，\(q=-2\).
因为 \(q\ne1\)，前 \(n\) 项和为
\[
a_{1}+a_{2}+\cdots+a_{n}
=a_{1}\frac{1-q^{n}}{1-q}
=\frac{1}{2}\cdot\frac{1-(-2)^{n}}{3}
=\frac{1-(-2)^{n}}{6}
\]
\end{solution}

\begin{problem}
已知函数 \(f(x) = \begin{cases}
\dfrac{2}{x}, & x \geqslant 2,\\
(x - 1)^3, & x < 2,
\end{cases}\)
若关于 \(x\) 的方程 \(f(x) = k\) 有两个不同的实根，则实数 \(k\) 的取值范围是\fillinblank{}.
\end{problem}

\begin{answer}
\((0,1)\)
\end{answer}

\begin{solution}
当 \(x<2\) 时，\(f(x)=(x-1)^{3}\)，其值域为
\[
(-\infty,1)
\]
且对每个 \(k<1\) 都有唯一的根
\[
x=1+\sqrt[3]{k}<2
\]
当 \(x\geqslant2\) 时，\(f(x)=\frac{2}{x}\)，其值域为
\[
(0,1]
\]
且当 \(0<k\leqslant1\) 时有唯一的根 \(x=\frac{2}{k}\geqslant2\).
因此，方程 \(f(x)=k\) 有两个不同实根，当且仅当两个分段都能
提供根且 \(k<1\)，即
\[
0<k<1
\]
\end{solution}

\begin{problem}
设 \(A(0,0)\)，\(B(4,0)\)，\(C(t + 4,3)\)，\(D(t,3) \)（\(t \in \R\)）. 记 \(N(t)\) 为平行四边形 \(ABCD\) 内部（不含边界）的整点的个数. 其中整点是指横，纵坐标都是整数的点，则 \(N(0) =\)\fillinblank{}；\(N(t)\) 的所有可能取值为\fillinblank{}.
\end{problem}

\begin{answer}
\(6\)；\(6\)，\(7\)，\(8\)
\end{answer}

\begin{solution}
平行四边形内的点可唯一表示为
\(s(4,0)+r(t,3)\)，\(0<s<1\)，\(0<r<1\).
若该点为整点，则其纵坐标 \(3r\) 为整数，所以
\[
r=\frac{1}{3}\quad\text{或}\quad r=\frac{2}{3}
\]
当 \(r=\frac{1}{3}\) 时，横坐标满足
\[
\frac{t}{3}<x<4+\frac{t}{3};
\]
当 \(r=\frac{2}{3}\) 时，横坐标满足
\[
\frac{2t}{3}<x<4+\frac{2t}{3}
\]
每个开区间长度均为 \(4\). 长度为 \(4\) 的开区间中，若左端点为
整数则含有 \(3\) 个整数，否则含有 \(4\) 个整数.

当 \(t=0\) 时，两行分别为 \(0<x<4\)，每行有 \(3\) 个整数，
所以 \(N(0)=6\).

第一行左端点为整数当且仅当 \(t\in3\Z\)，第二行左端点为
整数当且仅当 \(t\in\frac{3}{2}\Z\). 两行各有 \(3\) 或
\(4\) 个整数，且两行同时为 \(3\) 个、恰有一行为 \(3\) 个、或
同时为 \(4\) 个均可取到，例如分别取 \(t=0\)，\(\frac{3}{2}\)，\(1\).
故 \(N(t)\) 的所有可能取值为 \(6\)，\(7\)，\(8\).
\end{solution}

\section{解答题}

\begin{problem}
已知函数 \( f(x) = 4\cos x\sin \left(x + \frac{\pi}{6}\right) - 1 \).
\begin{enumerate}
\item 求 \( f(x) \) 的最小正周期.
\item 求 \( f(x) \) 在区间 \( \left[-\frac{\pi}{6}, \frac{\pi}{4}\right] \) 上的最大值和最小值.
\end{enumerate}
\end{problem}

\begin{solution}
\begin{enumerate}
\item 利用积化和差公式，
\[
f(x)=4\cos x\sin\left(x+\frac{\pi}{6}\right)-1
=2\left[\sin\left(2x+\frac{\pi}{6}\right)+\sin\frac{\pi}{6}\right]-1
=2\sin\left(2x+\frac{\pi}{6}\right)
\]
因此 \(f(x)\) 的周期为 \(\pi\). 因为正弦函数的最小正周期为
\(2\pi\)，所以 \(f(x)\) 的最小正周期为 \(\pi\).
\item 令
\[
t=2x+\frac{\pi}{6}
\]
当 \(x\in\left[-\frac{\pi}{6},\frac{\pi}{4}\right]\) 时，
\[
t\in\left[-\frac{\pi}{6},\frac{2\pi}{3}\right]
\]
在此区间内，\(\sin t\) 的最大值为 \(1\)，在 \(t=\frac{\pi}{2}\)
时取得；最小值为 \(-\frac{1}{2}\)，在 \(t=-\frac{\pi}{6}\)
时取得. 因而
\(\max f(x)=2\)，\(x=\frac{\pi}{6};\)，\(\min f(x)=-1\)，\(x=-\frac{\pi}{6}\).
\end{enumerate}
\end{solution}

\begin{problem}
以下茎叶图记录了甲，乙两组各四名同学的植树棵数. 乙组记录中有一个数据模糊，无法确认，在图中以 \(X\) 表示.

\begin{center}
\begin{tabular}{r|c|l}
甲组 & & 乙组 \\
\hline
\(9\quad 9\) & 0 & \(X\)，\(8\)，\(9\) \\
\(1\quad 1\) & 1 & \(0\)
\end{tabular}
\end{center}

\begin{enumerate}
\item 如果 \(X = 8\)，求乙组同学植树棵数的平均数和方差；
\item 如果 \(X = 9\)，分别从甲，乙两组中随机选取一名同学，求这两名同学的植树总棵数为 19 的概率.
\end{enumerate}
（注：方差 \( s^2 = \frac{1}{n}\left[(x_1 - \overline{x})^2 + (x_2 - \overline{x})^2 + \ldots + (x_n - \overline{x})^2\right] \)，其中 \(\overline{x}\) 为 \(x_1\)，\(x_2\)，\(\cdots\)，\(x_n\) 的平均数）

\end{problem}

\begin{solution}
\begin{enumerate}
\item 当 \(X=8\) 时，乙组数据为 \(8\)，\(8\)，\(9\)，\(10\)，所以平均数为
\[
\overline{x}=\frac{8+8+9+10}{4}=\frac{35}{4}
\]
由题给方差定义，
\[
s^{2}=\frac{1}{4}\left[
\left(8-\frac{35}{4}\right)^{2}
+\left(8-\frac{35}{4}\right)^{2}
+\left(9-\frac{35}{4}\right)^{2}
+\left(10-\frac{35}{4}\right)^{2}\right]
=\frac{1}{4}\cdot\frac{9+9+1+25}{16}
=\frac{11}{16}
\]
\item 当 \(X=9\) 时，甲组数据为 \(9\)，\(9\)，\(11\)，\(11\)，乙组数据为
\(9\)，\(8\)，\(9\)，\(10\). 从两组各选一人共有 \(4\times4=16\) 种等可能结果.
总数为 \(19\) 的情况是甲组取 \(9\)、乙组取 \(10\)，有 \(2\) 种，
或甲组取 \(11\)、乙组取 \(8\)，有 \(2\) 种. 因此有利结果有
\(4\) 种，所求概率为
\[
\frac{4}{16}=\frac{1}{4}
\]
\end{enumerate}
\end{solution}

\begin{problem}
如图，在四面体 \(PABC\) 中，\(PC \perp AB\)，\(PA \perp BC\)，点 \(D\)，\(E\)，\(F\)，\(G\) 分别是棱 \(AP\)，\(AC\)，\(BC\)，\(PB\) 的中点.

\begin{enumerate}
\item 求证：\(DE \parallel\) 平面 \(BCP\)；
\item 求证：四边形 \(DEFG\) 为矩形；
\item 是否存在点 \(Q\)，到四面体 \(PABC\) 六条棱的中点的距离相等？说明理由.
\end{enumerate}

\texfigure[max width=0.48\linewidth]{img_repaint/2011/beijing_liberal/q17_fig1.tex}
\end{problem}

\begin{solution}
\begin{enumerate}
\item 在三角形 \(APC\) 中，\(D\)，\(E\) 分别为 \(AP\)，\(AC\) 的中点，
所以由中位线定理
\[
DE\parallel PC
\]
又 \(PC\) 在平面 \(BCP\) 内，且四面体非退化时平面 \(APC\) 与平面
\(BCP\) 的交线为 \(PC\)，而 \(DE\subset\) 平面 \(APC\) 且
\(DE\ne PC\)，所以 \(DE\) 不在平面 \(BCP\) 内，
所以 \(DE\parallel\) 平面 \(BCP\).
\item 在三角形 \(APC\) 和三角形 \(PBC\) 中分别使用中位线定理，得
\(DE\parallel PC\parallel FG\)，\(EF\parallel AB\parallel DG\).
所以四边形 \(DEFG\) 的两组对边分别平行，是平行四边形.
由已知 \(PC\perp AB\)，且 \(DE\parallel PC\)、\(DG\parallel AB\)，
可得 \(DE\perp DG\). 平行四边形有一个直角，所以 \(DEFG\) 为矩形.
\item 设 \(M\)，\(N\) 分别为 \(PC\)，\(AB\) 的中点，设 \(Q\) 为矩形
\(DEFG\) 两条对角线的交点. 矩形的中心到四个顶点距离相等，
故
\[
QD=QE=QF=QG
\]
在四边形 \(MENG\) 中，\(M\)，\(E\) 分别为 \(PC\)，\(AC\) 的中点，
故 \(ME\parallel PA\)；\(E\)，\(N\) 分别为 \(AC\)，\(AB\) 的中点，
故 \(EN\parallel CB\). 同理 \(NG\parallel PA\)、\(GM\parallel CB\).
因此 \(MENG\) 为平行四边形. 又 \(PA\perp BC\)，所以它是矩形.
其对角线 \(MN\) 与 \(EG\) 互相平分，而 \(Q\) 是 \(EG\) 的中点，
故 \(Q\) 也是 \(MN\) 的中点，从而
\[
QM=QN=QE=QG
\]
结合前式，\(Q\) 到六条棱的中点 \(D\)，\(E\)，\(F\)，\(G\)，\(M\)，\(N\) 的距离均相等，
所以这样的点 \(Q\) 存在.
\end{enumerate}
\end{solution}

\begin{problem}
已知函数 \( f(x) = (x - k)\e^x \).
\begin{enumerate}
\item 求 \( f(x) \) 的单调区间；
\item 求 \( f(x) \) 在区间 \( [0,1] \) 上的最小值.
\end{enumerate}
\end{problem}

\begin{solution}
\begin{enumerate}
\item 求导得
\[
f'(x)=\e^x+(x-k)\e^x=(x-k+1)\e^x
\]
由于 \(\e^x>0\)，所以 \(f'(x)\) 的符号由 \(x-k+1\) 决定.
当 \(x<k-1\) 时 \(f'(x)<0\)，当 \(x>k-1\) 时 \(f'(x)>0\).
因此 \(f(x)\) 在 \((-\infty,k-1)\) 上单调递减，在
\((k-1,+\infty)\) 上单调递增.
\item 在区间 \([0,1]\) 上按临界点 \(k-1\) 的位置分类.
若 \(k\leqslant1\)，则 \(k-1\leqslant0\)，函数在 \([0,1]\) 上递增，
最小值为
\[
f(0)=-k
\]
若 \(1<k<2\)，则 \(0<k-1<1\)，最小值在临界点取得：
\[
f(k-1)=((k-1)-k)\e^{k-1}=-\e^{k-1}
\]
若 \(k\geqslant2\)，则 \(k-1\geqslant1\)，函数在 \([0,1]\) 上递减，
最小值为
\[
f(1)=(1-k)\e
\]
综上，在区间 \([0,1]\) 上的最小值为
\[
\begin{cases}
-k, & k\leqslant 1,\\
-\e^{k-1}, & 1<k<2,\\
(1-k)\e, & k\geqslant 2.
\end{cases}
\]
\end{enumerate}
\end{solution}

\begin{problem}
已知椭圆 \(G: \frac{x^2}{a^{2}} + \frac{y^2}{b^{2}} = 1 \) （\(a > b > 0\)）的离心率为 \(\frac{\sqrt{6}}{3}\)，右焦点为 \((2\sqrt{2}, 0)\)，斜率为 1 的直线 \(l\) 与椭圆 \(G\) 交于 \(A\)，\(B\) 两点. 以 \(AB\) 为底边作等腰三角形，顶点为 \(P(-3, 2)\).
\begin{enumerate}
\item 求椭圆 \(G\) 的方程；
\item 求 \( \triangle PAB \) 的面积.
\end{enumerate}
\end{problem}

\begin{solution}
\begin{enumerate}
\item 设椭圆半焦距为 \(c\). 由右焦点坐标，
\[
c=2\sqrt{2}
\]
又由离心率 \(e=\frac{c}{a}=\frac{\sqrt{6}}{3}\)，得
\[
a=\frac{c}{e}
=\frac{2\sqrt{2}}{\sqrt{6}/3}=2\sqrt{3}
\]
所以 \(a^{2}=12\). 由 \(c^{2}=a^{2}-b^{2}\)，得
\[
b^{2}=12-8=4
\]
故椭圆方程为
\[
\frac{x^{2}}{12}+\frac{y^{2}}{4}=1
\]
\item 设直线 \(l\) 的方程为 \(y=x+m\). 代入椭圆方程，得
\[
\frac{x^{2}}{12}+\frac{(x+m)^{2}}{4}=1
\]
即
\[
4x^{2}+6mx+3m^{2}-12=0
\]
两交点 \(A\)，\(B\) 的横坐标之和为 \(-\frac{3m}{2}\)，故弦 \(AB\)
的中点为
\[
E\left(-\frac{3m}{4},\frac{m}{4}\right)
\]
因为 \(PAB\) 是以 \(AB\) 为底边的等腰三角形，所以 \(PE\perp AB\).
直线 \(AB\) 的斜率为 \(1\)，故 \(PE\) 的斜率为 \(-1\). 当
\(m\ne4\) 时，
\[
\frac{\frac{m}{4}-2}{-\frac{3m}{4}+3}=-1
\]
解得 \(m=2\). 当 \(m=4\) 时，\(E=(-3,1)\)，直线 \(PE\)
为竖直线，也不满足与斜率为 \(1\) 的直线垂直，故仍只有 \(m=2\).

取 \(m=2\)，交点横坐标满足
\[
4x^{2}+12x=0
\] 所以 \(x=-3\)，\(0\)，从而
\(A=(-3,-1)\)，\(B=(0,2)\).
因此
\(AB=3\sqrt{2}\)，\(\operatorname{dist}(P,AB) =\frac{|-3-2+2|}{\sqrt{2}} =\frac{3}{\sqrt{2}}\).
故
\[
[\triangle PAB]
=\frac{1}{2}\cdot3\sqrt{2}\cdot\frac{3}{\sqrt{2}}
=\frac{9}{2}
\]
\end{enumerate}
\end{solution}

\begin{problem}
若数列 \(A_{n}\)：\(a_{1}\)，\(a_{2}\)，\(\cdots\)，\(a_{n}\)（\(n \geqslant 2\)）满足 \(\left|a_{k+1} - a_{k}\right| = 1\)（\(k = 1\)，\(2\)，\(\cdots\)，\(n-1\)），则称 \(A_{n}\) 为 \(E\) 数列. 记 \(S(A_{n}) = a_{1} + a_{2} + \cdots + a_{n}\).
\begin{enumerate}
\item 写出一个 \(E\) 数列 \(A_{n}\) 满足 \(a_{1}=a_{3}=0\)；
\item 若 \(a_1 = 12\)，\(n = 2000\)，证明：\(E\) 数列 \(A_n\) 是递增数列的充要条件是 \(a_n = 2011\)；
\item 在 \(a_1 = 4\) 的 \(E\) 数列 \(A_n\) 中，求使得 \(S(A_n) = 0\) 成立的 \(n\) 的最小值.
\end{enumerate}
\end{problem}

\begin{solution}
\begin{enumerate}
\item 例如取
\[
A_5=(0,1,0,1,0)
\]
相邻两项之差的绝对值均为 \(1\)，且 \(a_1=a_3=0\)，所以这是一个
满足条件的 \(E\) 数列.
\item 先证必要性. 若 \(A_n\) 是递增数列，则对每个
\(j=1\)，\(2\)，\(\ldots\)，\(1999\)，有
\[
a_{j+1}-a_j>0
\]
又因 \(\left|a_{j+1}-a_j\right|=1\)，所以
\[
a_{j+1}-a_j=1
\]
逐项相加得
\[
a_{2000}=a_1+\sum_{j=1}^{1999}(a_{j+1}-a_j)
=12+1999=2011
\]

再证充分性. 若 \(a_{2000}=2011\)，则
\[
\sum_{j=1}^{1999}(a_{j+1}-a_j)=a_{2000}-a_1=1999
\]
每一项 \(a_{j+1}-a_j\) 只能取 \(1\) 或 \(-1\)，特别地都不超过
\(1\). 共有 \(1999\) 项且它们的和等于 \(1999\)，因此每一项都必须
等于 \(1\). 于是 \(a_{j+1}>a_j\) 对所有 \(j\) 成立，\(A_n\) 是递增
数列. 故 \(A_n\) 是递增数列的充要条件为 \(a_n=a_{2000}=2011\).
\item 由 \(\left|a_{j+1}-a_j\right|=1\)，有
\[
a_j\geqslant a_{j-1}-1
\]
由 \(a_1=4\) 归纳可得
\[
a_j\geqslant4-(j-1)=5-j
\]
因此当 \(2\leqslant n\leqslant8\) 时，
\[
S(A_n)=\sum_{j=1}^{n}a_j
\geqslant\sum_{j=1}^{n}(5-j)
=\frac{n(9-n)}{2}>0
\]
不可能有 \(S(A_n)=0\). 取
\[
A_9=(4,3,2,1,0,-1,-2,-3,-4)
\]
它是 \(E\) 数列，且
\[
S(A_9)=4+3+2+1+0-1-2-3-4=0
\]
所以使 \(S(A_n)=0\) 成立的 \(n\) 的最小值为 \(9\).
\end{enumerate}
\end{solution}
